%%%%%%%%%%%%%%%%%%%%%%%%%%%%%%%%%%%%%%%%%%%%%%%%%%%%%%%%%%%%%%%%%%%%%%%%%%%%%%%
%% Capítulo 2 --- Fundamentação e revisão da literatura
%%
%% Reordenado pela ordem de precedência da cadeia, na MESMA sequência do
%% Capítulo 4: aquisição -> normalização -> estimação -> implantação.
%% Os relatórios de origem organizavam a revisão por "eixos" numerados 1, 3, 4 e
%% 5, abrindo pelo detector. A ordem aqui é a da cadeia, para que o capítulo de
%% revisão e o de resultados possam ser lidos em paralelo.
%%
%% A ausência de um eixo na numeração original era deliberada e está declarada
%% na Seção 2.1: as referências que o sustentariam não passaram na verificação.
%%%%%%%%%%%%%%%%%%%%%%%%%%%%%%%%%%%%%%%%%%%%%%%%%%%%%%%%%%%%%%%%%%%%%%%%%%%%%%%

\makeatletter
\@ifundefined{toprule}{%
  \newcommand{\toprule}{\hline}%
  \newcommand{\midrule}{\hline}%
  \newcommand{\bottomrule}{\hline}%
}{}
\@ifundefined{emandamento}{%
  \@ifundefined{textcolor}%
    {\newcommand{\emandamento}[1]{#1}}%
    {\newcommand{\emandamento}[1]{\textcolor{red}{#1}}}%
}{}
\makeatother

\chapter{Fundamentação e revisão da literatura}
\label{cap:revisao}

\section{Método da revisão}
\label{sec:rev-metodo}

A revisão foi conduzida por confronto, e não por levantamento: cada seção parte de uma
escolha de projeto ou de um resultado obtido e pergunta o que a literatura publicada diz
sobre ele --- se o confirma, se o contradiz ou se nada diz.

Todas as referências foram conferidas em fonte primária: depósito do editor, texto integral
do artigo, amostra oficial da norma ou folha de dados do fabricante. Onde o texto integral
não pôde ser acessado, isso consta na própria referência, e o achado correspondente é
tratado como de menor peso. Essa verificação não é formalidade: duas auditorias sucessivas
encontraram defeito em parcela substancial das referências inicialmente levantadas ---
primeiro autor trocado, iniciais incorretas colapsando autores distintos, norma retirada
citada como vigente e, em três casos, achado atribuído a uma fonte que não o continha.

Um eixo temático levantado na fase inicial não integra esta revisão. Nenhuma das
referências que o sustentariam sobreviveu à verificação em fonte primária, e a decisão foi
não citar o que não foi conferido. O registro fica aqui porque uma revisão que apresenta
apenas os eixos sobreviventes esconde o método que os selecionou.

A organização segue a ordem de precedência da cadeia de medição, a mesma do
Capítulo~\ref{cap:resultados}: primeiro a aquisição, que decide que informação existe;
depois a normalização, que impõe requisitos sobre o instrumento e sobre a estimativa; em
seguida a estimação, que decide o que se faz com a informação disponível; por fim a
implantação, que decide por quanto tempo o dispositivo continua operando.

\section{Aquisição: banda, ruído e detectabilidade de defeito de rolamento}
\label{sec:rev-aquisicao}

A tensão que organiza esta literatura é única: as frequências que se deseja medir em um
rolamento são baixas, mas o mecanismo físico que as torna mensuráveis é alto. O defeito
produz impactos periódicos na frequência característica da geometria, na ordem de dezenas
a centenas de hertz, e esses impactos excitam ressonâncias estruturais de alguns
quilohertz. A assinatura diagnóstica viaja modulada nessa portadora de alta frequência, e é
por demodulação de envelope sobre a banda de ressonância que ela é recuperada
\cite{randall2011,antoni2006,antoni2007}. O procedimento padrão localiza a banda ótima pela
díade que maximiza a impulsividade do sinal filtrado \cite{antoni2007}; um dispositivo cuja
frequência de Nyquist é de \SI{500}{\hertz} não tem díade a escolher.

Os estudos de adequação de sensor convergem para o limite superior de banda como a variável
determinante, e nenhum testa dispositivo próximo da faixa do Kaelix. A comparação de três
acelerômetros MEMS capacitivos contra um piezoelétrico de referência reprovou o de banda
mais estreita, de \SI{1500}{\hertz}, para monitoramento de condição \cite{albarbar2008}.
Os grupos que se propõem a diagnosticar rolamento com MEMS escolhem sensores de dezenas de
quilohertz \cite{landi2023,sharma2024}, e uma implementação de baixo custo publicada
resolve o problema por projeto, reservando um acelerômetro de banda estreita para
desbalanceamento e outro de \SI{40}{\kilo\hertz} para rolamento \cite{jakobsen2024}. Os
produtos comerciais medem, ao menos em um eixo, uma ordem de grandeza acima do teto de uma
IMU de consumo \cite{bentlynevada2025,erbessd2025,tractian2022}.

Há uma segunda objeção, anterior à banda e independente dela: o piso de ruído. A densidade
de ruído declarada em folha de dados para uma IMU de consumo, integrada sobre a banda
normativa, situa o sensor acima do ruído admissível para resolver a fronteira entre as
zonas normativas mais baixas. A literatura de MEMS levanta, de forma convergente, a
ausência de calibração rastreável desses dispositivos \cite{albarbar2008,rodrigues2021},
e a avaliação de sensibilidade diagnóstica de sensores digitais de catálogo chega à mesma
leitura \cite{fidali2024}. \emandamento{O orçamento de ruído usualmente atribuído à
literatura de fabricante não pôde ser confirmado em fonte primária \cite{looney2017}, e
fechar essa conta com densidade medida, e não declarada, permanece pendente.}

Um efeito adicional é pouco documentado e afeta diretamente dispositivos de consumo: a
ausência de filtro anti-\emph{aliasing} configurado abaixo da frequência de Nyquist. Sem
ele, conteúdo espectral acima da banda dobra para dentro da faixa de análise e torna-se
indistinguível de sinal legítimo. Não se localizou, na literatura revisada por pares de
diagnóstico de rolamento, descrição experimental do efeito desse artefato sobre descritores
estatísticos como curtose e fator de crista; ele aparece apenas em notas de aplicação de
fabricante, não conferidas nesta revisão.

\section{Normalização: o que a norma exige do instrumento e da estimativa}
\label{sec:rev-norma}

A avaliação de severidade em velocidade eficaz não é convenção arbitrária de unidades, mas
escolha justificada no próprio texto normativo \cite{iso208163_2022}. A norma vigente
estabelece as zonas de severidade para máquinas acima de \SI{15}{\kilo\watt} com rotação
entre \num{120} e \num{30000}~r/min, e exige do instrumento resposta plana em faixa de
\emph{ao menos} \SIrange{10}{1000}{\hertz}, com limite inferior de até \SI{2}{\hertz} em
máquinas próximas ou abaixo de \num{600}~r/min.

A distinção entre requisito de instrumento e janela de análise é decisiva e frequentemente
perdida. A exigência de resposta plana recai sobre o instrumento; é isso que torna o
truncamento de banda uma violação de requisito, e não uma escolha de processamento
reversível. A edição anterior da norma delega esse ``plana'' a uma norma de instrumentação
que o torna verificável, com tolerância de banda passante de aproximadamente
\SI{10}{\percent} entre \SI{12,6}{\hertz} e \SI{1000}{\hertz} e limitação de banda por par
de Butterworth de terceira ordem \cite{iso2954_2012}.

\emandamento{A edição de 2009 foi retirada e substituída pela de 2022
\cite{iso108163_2009,iso208163_2022}. Os valores de zona são idênticos nas duas, de modo
que a substituição afeta a rastreabilidade da citação e não os números.} A edição anterior
é mantida na bibliografia por ser a que a prática industrial e parte da literatura ainda
adotam, e por ser nela que está a delegação explícita à norma de instrumentação.

A conversão de aceleração para velocidade é obrigação normativa, e é também um problema
numérico com solução conhecida. A integração no domínio do tempo acumula qualquer
componente contínua presente no sinal, e o erro resultante cresce sem limite; a integração
no domínio da frequência, restrita à máscara normativa, é o método estabelecido, e a
comparação quantitativa entre os dois está publicada \cite{brandt2014}. O que não se
localizou publicado é a aplicação desse método ao indicador da norma em um dispositivo
truncado --- isto é, quanto um instrumento limitado a \SI{500}{\hertz} subestima o
indicador normativo frente à banda exigida.

\section{Estimação: detector, limiar, métrica e protocolo}
\label{sec:rev-estimacao}

Esta literatura desloca a pergunta de engenharia, mas o deslocamento é condicional, e a
condição precisa ficar à vista: \emph{dada banda de medição suficiente}, o que decide o
desempenho não é qual detector se escolhe, e sim como o limiar é calibrado, com que métrica
o resultado é lido e quanto do orçamento do microcontrolador o modelo consome. Sem banda,
nada disso importa --- nenhum corte de escore transforma contraste ausente em detecção.

\subsection{A escolha do detector não tem consenso}

Dois comparativos de grande porte discordam, e a discordância é informativa porque medem
coisas diferentes. O maior em cobertura de paradigmas avalia \num{30} algoritmos sobre
\num{57} conjuntos de dados e conclui que \emph{nenhum} método não supervisionado é
estatisticamente superior aos demais \cite{han2022}. A maior comparação restrita a métodos
não supervisionados, com \num{33} algoritmos sobre \num{52} conjuntos tabulares reais,
encontra vantagem para o Extended Isolation Forest, com a ressalva de que os dados se
separam em dois regimes \cite{bouman2024}. Não há contradição, e sim escopos distintos ---
mas há consequência: a evidência disponível recomenda uma variante que este trabalho não
adota.

A razão para permanecer no algoritmo canônico \cite{liu2008} é de custo de implementação
embarcada, e é verificável em estrutura: um nó da variante estendida armazena um vetor
normal com tantas componentes quantos forem os descritores, e a inferência passa de uma
comparação para um produto interno por nó. \emandamento{É argumento de projeto, e não
resultado de literatura; enquanto o orçamento embarcado não for medido no dispositivo, a
escolha permanece limitação de escopo declarada, não conclusão.}

O suporte específico para máquinas rotativas é, em contrapartida, o melhor que o algoritmo
canônico tem: entre \num{11} detectores não supervisionados, atinge F1 de \num{97,19} em
rolamento \emph{run-to-failure} e \num{97,20} em faltas mecânicas de bancada
\cite{brito2022}. Os dois primeiros conjuntos são adquiridos a \SI{20}{\kilo\hertz}, banda
que uma IMU de consumo não possui.

\subsection{O limiar padrão não é calibração}

No artigo fundador, o escore tem \num{0,5} como \emph{referência interpretativa}, e a
robustez demonstrada está nos parâmetros do \emph{ensemble}, não no corte \cite{liu2008}.
A implementação de referência explicita o mecanismo: o parâmetro de contaminação redefine
apenas um deslocamento fixo, nunca o escore subjacente \cite{scikitlearn2026}. Seguem daí
duas consequências mecânicas: o corte padrão é herdado de uma interpretação, e não estimado
nos dados; e toda métrica de ranqueamento, área sob a curva ROC inclusive, é invariante ao
parâmetro.

Os trabalhos que reportam bom desempenho abandonam o padrão --- fixando o escore
normalizado em \num{0,75} sob hipótese de treino limpo \cite{antonini2023}, ou derivando o
corte da razão de contaminação conhecida \cite{brito2022}. A alternativa principiada
dispensa hipótese sobre a distribuição do corpo dos dados e fixa o limiar a partir de um
risco alvo, que controla diretamente o número de falsos positivos \cite{siffer2017}.

O problema não é teórico. Três implementações independentes com o mesmo detector em nó
embarcado situam-se entre \num{22} e \SI{38}{\percent} de falso alarme sobre operação
normal \cite{chevtchenko2023,kolok2025}. Um dispositivo que marca um terço da operação
normal como anomalia não é utilizável em manutenção.

\subsection{A métrica e o protocolo de validação}

A área sob a curva ROC global não caracteriza o ponto de operação de um dispositivo de
campo, que precisa operar em falso alarme baixo. O benchmark de referência para o mesmo
regime de treino reporta área parcial restrita a essa faixa \cite{koizumi2020}, e a mesma
disciplina aparece em quem reporta PR-AUC além do F1 \cite{brito2022}.

Sobre o protocolo, a objeção é a mais séria e não foi identificada por este trabalho.
Janelas consecutivas de um mesmo ensaio compartilham a assinatura da montagem, e sua
divisão aleatória entre treino e teste permite que o modelo identifique o ensaio em vez da
falha. O fenômeno tem nome e quantificação: de \num{41} estudos revisados sobre uma base de
referência consolidada, \num{40} adotaram delineamento suscetível a ele \cite{varejao2025},
e a inflação medida por vazamento é da ordem de dezenas de pontos percentuais
\cite{vieira2026}. Trabalhos que revisam abordagens de classificação de falha de rolamento
chegam a leitura convergente \cite{abburi2023,smith2015}.

\section{Implantação: nós comparáveis, invólucro e energia}
\label{sec:rev-implantacao}

A literatura de nós sem fio para manutenção preditiva divide-se em duas populações que
raramente se citam. De um lado, produtos comerciais, com banda uma ordem de grandeza acima
da de uma IMU de consumo, caminho mecânico rígido até o elemento sensor e alimentação por
célula primária \cite{bentlynevada2025,erbessd2025,tractian2022}. De outro, nós acadêmicos
de baixo custo, que convergem para a mesma arquitetura estudada aqui: microcontrolador de
baixo consumo, IMU MEMS na ordem de \SI{1}{\kilo\hertz} e detecção não supervisionada
\cite{kolok2025,antonini2023,bisio2025}.

\subsection{O que a literatura de nós de baixo custo autoriza concluir}

O análogo mais próximo em hardware e em modo de falha reporta acurácia de
\SI{72,9}{\percent} $\pm$~\SI{2,4}{\percent}, com validação cruzada de cinco dobras sobre
cerca de \num{900} amostras e \emph{um único} modo de falha \cite{kolok2025}. O dado valida
a arquitetura como publicável. O que ele não faz é fixar teto de desempenho da classe: a
amostra é de um estudo único com um modo de falha, e o trabalho encerra em valor eficaz de
aceleração, sem integração para velocidade e sem critério normativo.

A comparação mais direta é embarcada e mede três detectores no mesmo nó e no mesmo
conjunto \cite{chevtchenko2023}. No mesmo hardware e no mesmo dado, um método baseado em
vizinhança domina o algoritmo adotado aqui tanto em especificidade quanto em latência,
sendo cerca de \num{26} vezes mais rápido. \emandamento{A resposta declarada é de custo de
memória --- o método de vizinhança exige manter o conjunto de referência, enquanto a
floresta armazena apenas limiares e índices --- mas o orçamento embarcado ainda não foi
medido, e a permanência no detector adotado é hipótese de projeto, não conclusão.}

O mesmo trabalho traz um achado que é restrição de implantação, e não detalhe: um motor com
cerca de um ano de uso tem assinatura normal distinta da de motores novos
\cite{chevtchenko2023}. Um modelo único com limiar fixo herda essa deriva do que conta como
normal entre máquinas e ao longo da vida --- cenário em que a calibração por risco alvo
deixa de ser refinamento e passa a ser requisito.

A viabilidade computacional, em contrapartida, está estabelecida: o detector treina e
executa dentro de um microcontrolador da classe \cite{antonini2023}, existe variante
desenhada para reduzir seu custo em TinyML \cite{barbariol2022}, e dispositivos da mesma
família comportam \emph{pipelines} consideravelmente mais caros \cite{uselis2025,kumar2026}.

\subsection{Invólucro e caminho mecânico}

É na física do dispositivo que a literatura contraria o projeto mais diretamente. A prática
consolidada em acelerometria exige que a frequência natural do conjunto montado seja
suficientemente superior à maior frequência de interesse, e os produtos comerciais obtêm
isso por base metálica aparafusada ou colada ao ativo. Nenhum dos nós de baixo custo
revisados publica a frequência natural do conjunto invólucro-sensor-fixação: um deles
reporta apenas autorressonâncias de placa e estrutura entre \SIrange{6}{9}{\kilo\hertz}
\cite{jakobsen2024}, outro adia explicitamente o invólucro \cite{antonini2023}, e para um
terceiro o dado não foi localizado na parte acessível do texto \cite{bisio2025}.

\subsection{Energia e margem térmica}

Na energia, a literatura contraria o projeto com clareza. O campo adota célula primária, e
portanto nunca carrega nada \cite{nabavi2025}; a revisão térmica disponível estabelece o
mecanismo de degradação por temperatura, mas não o limite de carga \cite{ma2018}. O limite
de carga de uma célula secundária vem da folha de dados e das normas de segurança
aplicáveis, não da literatura de manutenção preditiva. \emandamento{O orçamento térmico
completo --- carcaça quente, invólucro fechado e célula carregando --- não foi localizado
publicado, e é ele que decide se a escolha de bateria sobrevive.}

\section{Síntese}
\label{sec:rev-sintese}

Lidos em conjunto, os eixos não formam uma lista de pontos fortes e fracos: formam uma
cadeia com ordem de precedência. Primeiro a aquisição --- piso de ruído, banda e caminho
mecânico ---, e nenhuma decisão posterior recupera informação que o sensor não coletou.
Depois a estimação, que decide o que se faz com a informação disponível. Por último a
implantação, que decide por quanto tempo o dispositivo continua operando.

A conclusão desta revisão é que \textbf{a literatura é mais forte onde o dispositivo é mais
frágil, no primeiro elo, e mais frágil onde o dispositivo tem algo a dizer, no segundo.}

\subsection{O que a literatura sustenta}

O problema do limiar é real e não é preciosismo, e a alternativa principiada existe. A
disciplina de reportar a métrica no ponto de operação, e não apenas de forma agregada, tem
respaldo em benchmark de referência. A viabilidade computacional da arquitetura está
estabelecida, e a exigência de integrar aceleração para velocidade é normativa, com método
correto estabelecido. A forma do limite térmico --- margem imposta pela célula e não pelo
invólucro --- tem precedente comercial direto.

\subsection{O que a literatura contesta}

As objeções são de peso, e nenhuma se resolve com ajuste de \emph{software}. A banda:
nenhum estudo de adequação revisado testou sensor com limite superior próximo ao desta
classe, e o mecanismo físico da assinatura de rolamento não admite contorno. O piso de
ruído: a objeção é anterior à banda e não depende de nenhuma escolha de processamento. O
detector: a maior comparação restrita a métodos não supervisionados recomenda outra
variante, e no mesmo hardware embarcado um método de vizinhança mostrou-se mais específico
e mais rápido. \emandamento{Nenhuma dessas três objeções foi respondida por medição neste
trabalho; as respostas oferecidas são critérios de projeto declarados, e permanecem
hipóteses até que o orçamento embarcado seja medido.}

\section{Lacunas identificadas}
\label{sec:rev-lacunas}

Cada lacuna abaixo declara o que foi procurado, onde, e o que não foi encontrado. Duas
lacunas enunciadas na fase inicial caíram na verificação e estão registradas como caídas,
porque uma revisão que apresenta apenas as lacunas sobreviventes esconde o método.

\textbf{Caíram na verificação.} A afirmação de que não existiria falso alarme publicado
para o detector em nó embarcado é falsa: dois trabalhos o reportam ou permitem derivá-lo
\cite{chevtchenko2023,kolok2025}. A afirmação de que ninguém quantificaria a fração de
operação normal marcada pela política de limiar padrão também é falsa, pelo mesmo motivo.

\begin{enumerate}
  \item \textbf{Efeito isolado do parâmetro de contaminação.} Não se localizou trabalho que
        variasse o parâmetro sobre a \emph{mesma} base e o \emph{mesmo} modelo reportando a
        fração de operação normal marcada. O experimento que falta é barato e definitivo.
  \item \textbf{Falso alarme de detector treinado no próprio dispositivo.} O único trabalho
        revisado que treina \emph{on-device} é justamente o que declara a avaliação de
        acurácia fora de escopo \cite{antonini2023} --- e esse é o modo de operação exigido
        de uma máquina sem histórico rotulado.
  \item \textbf{Contraste diagnóstico como função contínua da banda.} A literatura discute
        o efeito da banda qualitativamente, e o único estudo revisado que faz varredura de
        taxa de amostragem o faz para acurácia de classificador supervisionado, não para
        contraste do indicador \cite{elbouharrouti2024}.
  \item \textbf{Artefato de \emph{aliasing} sobre descritores estatísticos.} Não localizado
        em literatura revisada por pares; aparece apenas em notas de aplicação de
        fabricante.
  \item \textbf{Custo do truncamento sobre o indicador normativo.} O erro dos métodos de
        integração está quantificado \cite{brandt2014}, mas não a aplicação ao indicador da
        norma em dispositivo truncado.
  \item \textbf{Detectabilidade limitada por piso de ruído.} Não se localizou análise de
        defeito mínimo detectável em função da densidade de ruído para IMUs de consumo.
  \item \textbf{Frequência natural do conjunto montado.} Nenhum nó de baixo custo revisado
        a publica.
  \item \textbf{Orçamento térmico de carga de célula secundária} em nó montado em carcaça
        de motor.
  \item \textbf{Separar efeito de banda de efeito de montagem.} Não se localizou benchmark
        aberto com sinais gravados \emph{simultaneamente} por MEMS de banda limitada e por
        piezoelétrico no mesmo ponto. Sem isso, os dois efeitos permanecem confundidos em
        qualquer medição de campo, inclusive nas deste projeto.
\end{enumerate}

As lacunas 1, 3 e 5 são as que o Capítulo~\ref{cap:resultados} aborda diretamente. As
demais permanecem abertas e reaparecem como trabalho futuro no
Capítulo~\ref{cap:conclusoes}.
